% !TEX root = master.tex
\chapter{Verzeichnisstruktur (Backend / Webserver)}
\label{ch:b-verzeichnis}

Wie Projekt~A wird auch SunShine~Tours vollst\"andig in das Webserver-Wurzelverzeichnis kopiert.
Die folgende \"Ubersicht zeigt die Struktur:

\begin{lstlisting}[numbers=none]
sunshine-tours/
|-- index.php              Startseite
|-- about.php              Ueber uns / Team
|-- tours.php              Touren und Kundenstimmen
|-- booking.php            Buchungsformular
|-- process_booking.php    Verarbeitung + Speicherung der Buchung
|-- confirmation.php       Buchungsbestaetigung
|-- legal.php              Impressum / Datenschutz / AGB
|-- formcheck.php          GET-/POST-Testseite
|-- visitor_count.txt      Zaehlerstand Besucherzaehler (beschreibbar)
|
|-- includes/
|   |-- header.php         Kopfbereich + Navigation + Theme-Schalter
|   |-- footer.php         Fussbereich + News-Ticker + Besucherzaehler
|
|-- js/
|   |-- booking.js         Echtzeit-Kosten, dynamische Felder, Validierung
|   |-- main.js            Hell-/Dunkel-Modus, mobiles Menue
|
|-- css/
|   |-- style.css          Farbschema, Layout, responsive Regeln
|
|-- images/
|   |-- paris.jpg, bielefeld.jpg, chernobyl.jpg,
|   |-- gary.jpg, pyongyang.jpg     (KI-generierte Bilder)
|
|-- bookings/
    |-- ST-JJJJMMTT-XXXXXX.json     gespeicherte Buchungen (beschreibbar)
\end{lstlisting}

\noindent Tabelle~\ref{tab:b-verzeichnis} fasst die Aufgaben der Komponenten zusammen.

\begin{small}
\begin{longtable}{@{}L{4.4cm} L{10.0cm}@{}}
\caption{Komponenten von Projekt~B}
\label{tab:b-verzeichnis}\\
\toprule
\textbf{Datei / Verzeichnis} & \textbf{Aufgabe} \\
\midrule
\endfirsthead
\toprule
\textbf{Datei / Verzeichnis} & \textbf{Aufgabe} \\
\midrule
\endhead
\bottomrule
\endfoot
\bottomrule
\endlastfoot
\texttt{index.php} & Startseite mit Hero-Bereich und Kennzahlen. \\
\texttt{about.php} & Vorstellung des fiktiven Teams und der Werte. \\
\texttt{tours.php} & Tourangebote und Kundenbewertungen. \\
\texttt{booking.php} & Buchungsformular mit Echtzeit-Kostenanzeige. \\
\texttt{process\_booking.php} & Wertet die Buchung aus, berechnet den Preis, maskiert die Kartennummer und speichert die Buchung als JSON. \\
\texttt{confirmation.php} & Zeigt die Zusammenfassung der Buchung. \\
\texttt{legal.php} & Impressum, Datenschutz und AGB. \\
\texttt{formcheck.php} & Demonstration der GET-/POST-Verarbeitung. \\
\texttt{includes/} & Gemeinsamer Kopf- und Fu\ss{}bereich (Navigation, Theme-Schalter, News-Ticker, Besucherz\"ahler). \\
\texttt{js/} & Clientlogik: Buchungsberechnung (\texttt{booking.js}) und Oberfl\"achenfunktionen (\texttt{main.js}). \\
\texttt{css/} & Zentrales Stylesheet inkl.\ Hell-/Dunkel-Modus. \\
\texttt{images/} & KI-generierte St\"adtebilder. \\
\texttt{bookings/} & Ablageverzeichnis der gespeicherten Buchungen (zur Laufzeit beschrieben). \\
\texttt{visitor\_count.txt} & Persistenter Z\"ahlerstand des Besucherz\"ahlers. \\
\bottomrule
\end{longtable}
\end{small}
